% !TeX root = ../mathematical-methods1.tex
\chapter{Vectors}\label{chap:vectors}

This chapter follows the lecturer's vector algebra, including the \textbf{index-notation}. 
It provides a foundation for later vector
calculus. Bold letters replace the manuscript's underlined vector notation.
All coordinates refer to a right-handed orthonormal Cartesian system.

\section{A geometric approach to vectors}

Informally, a vector is an arrow with a direction and a length. More precisely,
a geometric vector is an equivalence class of directed line segments in
three-dimensional Euclidean space. The segment from $A$ to $B$ represents
$\overrightarrow{AB}$; $A$ is its initial point and $B$ its terminal point.
Two directed segments represent the same vector when they have the same length
and direction, even if their starting points differ. Parallel arrows pointing
in opposite directions do not represent the same nonzero vector.

Examples from mechanics are the position vector $\mathbf r=\overrightarrow{OP}$
of a point relative to a fixed origin $O$, and its velocity $\mathbf v$.
If the endpoints coincide, $\overrightarrow{AA}=\mathbf0$ is the zero vector,
which has zero magnitude and no distinguished direction.

\subsection{Addition, subtraction and scalar multiplication}
If $\mathbf v_1$ and $\mathbf v_2$ are successive displacements, their total
is $\mathbf v_1+\mathbf v_2$. Place the tail of the second arrow at the tip
of the first; the resultant joins the initial tail to the final tip.
The parallelogram construction shows
$\mathbf v_1+\mathbf v_2=\mathbf v_2+\mathbf v_1$.

\begin{center}
\begin{tikzpicture}[scale=0.9,>=Stealth]
 \draw[->,thick] (0,0)--(3,0) node[midway,below] {$\mathbf v_1$};
 \draw[->,thick] (0,0)--(1,2) node[midway,left] {$\mathbf v_2$};
 \draw[->,dashed] (3,0)--(4,2) node[midway,right] {$\mathbf v_2$};
 \draw[->,dashed] (1,2)--(4,2) node[midway,above] {$\mathbf v_1$};
 \draw[->,very thick,ink] (0,0)--(4,2)
   node[midway,above,sloped] {$\mathbf v_1+\mathbf v_2$};
\end{tikzpicture}
\end{center}

Subtraction reverses addition: if $\mathbf a=\mathbf b+\mathbf c$, then
$\mathbf c=\mathbf a-\mathbf b$. In particular, $-\mathbf v$ has the same
magnitude as $\mathbf v$ and the opposite direction, and
$\mathbf v+(-\mathbf v)=\mathbf0$.
Repeated addition gives $2\mathbf v=\mathbf v+\mathbf v$ and
$3\mathbf v=\mathbf v+\mathbf v+\mathbf v$.
For $\lambda\in\mathbb R$, the vector $\lambda\mathbf v$ has length
$|\lambda|\,|\mathbf v|$; for nonzero $\mathbf v$ it points in the same
direction when $\lambda>0$ and in the opposite direction when $\lambda<0$.
Also $0\mathbf v=\mathbf0$.

\subsection{Properties of the operations}
For vectors $\mathbf u,\mathbf v,\mathbf w$ and real numbers $\lambda,\mu$,
\begin{align*}
 \mathbf v+\mathbf w&=\mathbf w+\mathbf v,
   & (\mathbf u+\mathbf v)+\mathbf w&=\mathbf u+(\mathbf v+\mathbf w),\\
 \mathbf v+\mathbf0&=\mathbf v,
   & \mathbf v+(-\mathbf v)&=\mathbf0,\\
 (\lambda+\mu)\mathbf v&=\lambda\mathbf v+\mu\mathbf v,
   & \lambda(\mathbf v+\mathbf w)&=\lambda\mathbf v+\lambda\mathbf w,\\
 \lambda(\mu\mathbf v)&=(\lambda\mu)\mathbf v,
   & 1\mathbf v&=\mathbf v.
\end{align*}
These can be verified geometrically. Matrices, functions and polynomials also
obey these vector-space rules with suitable addition and scalar multiplication;
geometric arrows are one instance of a broader algebraic idea.

\begin{example}[Lecture Example 1: the midpoint theorem]
Let $ABC$ be a triangle, and let $D,E$ be the midpoints of $AB,BC$ respectively.
Show that $DE$ is parallel to $AC$ and has half its length.
\begin{align*}
 \overrightarrow{DE}
 &=\overrightarrow{DB}+\overrightarrow{BE}\\
 &=\frac12\overrightarrow{AB}+\frac12\overrightarrow{BC}\\
 &=\frac12(\overrightarrow{AB}+\overrightarrow{BC})
 =\frac12\overrightarrow{AC}.
\end{align*}
The positive scalar $1/2$ gives the claimed direction and length.
\end{example}

\section{Cartesian coordinates}

Fix an origin $O$ and perpendicular coordinate axes $x,y,z$ in a right-handed
orientation. Let $\mathbf i,\mathbf j,\mathbf k$ be the unit vectors along
these axes. A point $P=(x,y,z)$ has position vector
\[
 \mathbf r=\overrightarrow{OP}=x\mathbf i+y\mathbf j+z\mathbf k.
\]
Conversely, translate the tail of any vector to $O$ to write it uniquely as
\[
 \mathbf v=v_1\mathbf i+v_2\mathbf j+v_3\mathbf k
 =\begin{pmatrix}v_1\\v_2\\v_3\end{pmatrix}
 =\langle v_1,v_2,v_3\rangle.
\]
The numbers $v_1,v_2,v_3$ are its components. Vector operations become
\[
 \begin{pmatrix}u_1\\u_2\\u_3\end{pmatrix}
 +\begin{pmatrix}v_1\\v_2\\v_3\end{pmatrix}
 =\begin{pmatrix}u_1+v_1\\u_2+v_2\\u_3+v_3\end{pmatrix},
 \qquad
 \lambda\begin{pmatrix}v_1\\v_2\\v_3\end{pmatrix}
 =\begin{pmatrix}\lambda v_1\\\lambda v_2\\\lambda v_3\end{pmatrix}.
\]

\begin{example}[Lecture Example 2: coordinates of a midpoint]
Let $A=(a_1,a_2,a_3)$ and $B=(b_1,b_2,b_3)$, and let $M$ be their midpoint.
\begin{align*}
 \overrightarrow{OM}
 &=\overrightarrow{OA}+\frac12\overrightarrow{AB}\\
 &=\overrightarrow{OA}+\frac12(\overrightarrow{OB}-\overrightarrow{OA})\\
 &=\frac12(\overrightarrow{OA}+\overrightarrow{OB}).
\end{align*}
Thus
\[
 M=\left(\frac{a_1+b_1}{2},\frac{a_2+b_2}{2},\frac{a_3+b_3}{2}\right).
\]
\end{example}

\subsection{Length and unit vectors}
Pythagoras, applied first in a coordinate plane and then in space, gives
\[
 |\mathbf v|^2=v_1^2+v_2^2+v_3^2,
 \qquad |\mathbf v|=\sqrt{v_1^2+v_2^2+v_3^2}.
\]
Since $|\lambda\mathbf v|=|\lambda|\,|\mathbf v|$, every nonzero vector
can be normalised:
\[
 \widehat{\mathbf v}=\frac{\mathbf v}{|\mathbf v|},
 \qquad |\widehat{\mathbf v}|=
 \frac{|\mathbf v|}{|\mathbf v|}=1.
\]
A vector of length $1$ is called a unit vector. The zero vector cannot be normalised.

\section{Straight lines}

A line is specified by a point $P_0$ on it and a nonzero direction vector
$\mathbf v$. If $\mathbf r_0=\overrightarrow{OP_0}$, then its points satisfy
\[
 \boxed{\mathbf r=\mathbf r_0+t\mathbf v,\qquad t\in\mathbb R.}
\]
Writing $P_0=(x_0,y_0,z_0)$ and $\mathbf v=\langle v_1,v_2,v_3\rangle$
gives the parametric equations
\[
 x=x_0+tv_1,\qquad y=y_0+tv_2,\qquad z=z_0+tv_3.
\]
If all three direction components are nonzero, eliminate $t$ to obtain
\[
 \frac{x-x_0}{v_1}=\frac{y-y_0}{v_2}=\frac{z-z_0}{v_3}.
\]
If a component is zero, use the corresponding constant-coordinate equation
instead of dividing by zero.

\begin{example}[Lecture Example 3: a line through two points]
Find the parametric and Cartesian forms of the line through
$P_0=(1,2,3)$ and $Q_0=(2,1,-1)$.
Its direction is
\[
 \overrightarrow{P_0Q_0}=\overrightarrow{OQ_0}-\overrightarrow{OP_0}
 =\langle2,1,-1\rangle-\langle1,2,3\rangle=\langle1,-1,-4\rangle.
\]
Thus
\[
 \mathbf r=\langle1,2,3\rangle+t\langle1,-1,-4\rangle,
\]
so
\[
 x=1+t,\qquad y=2-t,\qquad z=3-4t,
 \qquad \frac{x-1}{1}=\frac{y-2}{-1}=\frac{z-3}{-4}.
\]
\end{example}

\section{The scalar or dot product}

The mechanical motivation is work done by a constant force $\mathbf F$
through a displacement $\mathbf d$:
$W=|\mathbf F||\mathbf d|\cos\theta$.
For nonzero vectors meeting at angle $\theta\in[0,\pi]$, define
\[
 \mathbf a\cdot\mathbf b=|\mathbf a||\mathbf b|\cos\theta.
\]
Set the product to zero if either vector is zero. Two nonzero vectors are
orthogonal exactly when $\theta=\pi/2$, equivalently
$\mathbf a\cdot\mathbf b=0$.

\subsection{Deriving the component formula}
The cosine rule gives
\[
 |\mathbf a-\mathbf b|^2=|\mathbf a|^2+|\mathbf b|^2
 -2|\mathbf a||\mathbf b|\cos\theta.
\]
Therefore
\begin{align*}
 \mathbf a\cdot\mathbf b
 &=\frac12\bigl(|\mathbf a|^2+|\mathbf b|^2-|\mathbf a-\mathbf b|^2\bigr)\\
 &=\frac12\sum_{i=1}^3\bigl(a_i^2+b_i^2-(a_i-b_i)^2\bigr)\\
 &=a_1b_1+a_2b_2+a_3b_3.
\end{align*}
In particular,
\[
 \boxed{\mathbf a\cdot\mathbf b
 =|\mathbf a||\mathbf b|\cos\theta
 =a_1b_1+a_2b_2+a_3b_3},
 \qquad
 \theta=\arccos\frac{\mathbf a\cdot\mathbf b}{|\mathbf a||\mathbf b|}.
\]
The angle formula requires both vectors to be nonzero.

\subsection{Properties}
For real $\lambda$,
\begin{align*}
 \mathbf a\cdot(\mathbf b+\mathbf c)&=\mathbf a\cdot\mathbf b+\mathbf a\cdot\mathbf c,\\
 (\lambda\mathbf a)\cdot\mathbf b&=\lambda(\mathbf a\cdot\mathbf b),\\
 \mathbf a\cdot\mathbf b&=\mathbf b\cdot\mathbf a.
\end{align*}
For example, the first identity follows by expanding components:
\begin{align*}
 \mathbf a\cdot(\mathbf b+\mathbf c)
 &=a_1(b_1+c_1)+a_2(b_2+c_2)+a_3(b_3+c_3)\\
 &=(a_1b_1+a_2b_2+a_3b_3)+(a_1c_1+a_2c_2+a_3c_3).
\end{align*}
The other two follow similarly.

\subsection{Equations of planes}
A plane through $P_0$ with a nonzero normal vector $\mathbf n$ consists of
points $P$ for which $\overrightarrow{P_0P}\perp\mathbf n$.
Writing their position vectors as $\mathbf r_0,\mathbf r$,
\begin{align*}
 (\mathbf r-\mathbf r_0)\cdot\mathbf n&=0,\\
 \mathbf r\cdot\mathbf n&=\mathbf r_0\cdot\mathbf n=d.
\end{align*}
If $\mathbf n=\langle n_1,n_2,n_3\rangle$, the \textbf{Cartesian equation} is
\[
 \boxed{n_1x+n_2y+n_3z=d}.
\]
Alternatively, with two nonparallel directions $\mathbf u,\mathbf v$
parallel to the plane,
\[
 \boxed{\mathbf r=\mathbf r_0+s\mathbf u+t\mathbf v,\qquad s,t\in\mathbb R.}
\]
The nonparallel condition ensures that this describes a plane rather than a line.

\subsection{Orthogonal projection}
For $\mathbf u\ne\mathbf0$, project $\mathbf v$ onto the direction of
$\mathbf u$. The signed component is $|\mathbf v|\cos\theta$, and hence
\begin{align*}
 \mathbf p=\operatorname{proj}_{\mathbf u}\mathbf v
 &=|\mathbf v|\cos\theta\frac{\mathbf u}{|\mathbf u|}\\
 &=\boxed{\frac{\mathbf v\cdot\mathbf u}{|\mathbf u|^2}\mathbf u}.
\end{align*}
Check that the remaining component is orthogonal to $\mathbf u$:
\[
 (\mathbf v-\mathbf p)\cdot\mathbf u
 =\mathbf v\cdot\mathbf u
 -\frac{\mathbf v\cdot\mathbf u}{|\mathbf u|^2}(\mathbf u\cdot\mathbf u)=0.
\]
\begin{center}
\begin{tikzpicture}[scale=0.85,>=Stealth]
 \draw[->] (0,0)--(5,0) node[right] {$\mathbf u$ direction};
 \draw[->,thick,ink] (0,0)--(3,2) node[above] {$\mathbf v$};
 \draw[->,thick] (0,0)--(3,0) node[midway,below] {$\mathbf p$};
 \draw[dashed] (3,0)--(3,2);
 \draw (2.75,0)--(2.75,0.25)--(3,0.25);
\end{tikzpicture}
\end{center}

\begin{example}[Lecture Example 4: angle and projection]
Let $\mathbf a=\langle1,2,1\rangle$ and
$\mathbf b=\langle-2,1,2\rangle$. Then
\[
 \mathbf a\cdot\mathbf b=-2+2+2=2,
 \quad |\mathbf a|=\sqrt6,\quad |\mathbf b|=3.
\]
Thus
\[
 \cos\theta=\frac{2}{3\sqrt6}=\frac{\sqrt6}{9},
 \qquad\theta=\arccos\left(\frac{\sqrt6}{9}\right).
\]
The projection of $\mathbf a$ onto $\mathbf b$ is
\[
 \operatorname{proj}_{\mathbf b}\mathbf a
 =\frac29\langle-2,1,2\rangle
 =\left\langle-\frac49,\frac29,\frac49\right\rangle.
\]
\end{example}

\section{The vector or cross product}

The motivation is torque. A force $\mathbf F$ applied at position $\mathbf r$
relative to a fixed point $O$ produces torque of magnitude
$|\mathbf r||\mathbf F|\sin\theta$. Its axis is perpendicular to both
$\mathbf r$ and $\mathbf F$, with the sense determined by a right-hand rule.
This motivates the torque vector $\boldsymbol\tau=\mathbf r\times\mathbf F$.
Viewed from the tip of $\boldsymbol\tau$ towards the origin, the induced
rotation is anticlockwise.

\begin{definition}[Cross product]
For nonzero, noncollinear vectors $\mathbf a,\mathbf b$, their cross product
$\mathbf c=\mathbf a\times\mathbf b$ is characterised by
\begin{enumerate}
 \item $|\mathbf c|=|\mathbf a||\mathbf b|\sin\theta$;
 \item $\mathbf c$ is perpendicular to both $\mathbf a$ and $\mathbf b$;
 \item $\mathbf a,\mathbf b,\mathbf c$ form a right-handed triple.
\end{enumerate}
If either vector is zero or the vectors are collinear, set their cross product
to zero. Nonzero collinear vectors are parallel or antiparallel.
\end{definition}

\subsection{Deriving the Cartesian formula}
Write $\mathbf a=\langle a_1,a_2,a_3\rangle$ and
$\mathbf b=\langle b_1,b_2,b_3\rangle$.
A perpendicular vector $\mathbf c=\langle x,y,z\rangle$ satisfies
\[
 a_1x+a_2y+a_3z=0,\qquad b_1x+b_2y+b_3z=0.
\]
If $D=a_1b_2-a_2b_1\ne0$, solving these equations gives
\[
 x=\frac{a_2b_3-a_3b_2}{D}z,\qquad
 y=\frac{a_3b_1-a_1b_3}{D}z.
\]
Thus the perpendicular direction is
\[
 \mathbf q=\langle a_2b_3-a_3b_2,\ a_3b_1-a_1b_3,\ a_1b_2-a_2b_1\rangle,
 \qquad\mathbf c=t\mathbf q.
\]
For noncollinear vectors at least one such minor is nonzero; if $D=0$,
choose another pair of coordinates. The component expression for $\mathbf q$
remains the same. Expanding and collecting terms gives
\begin{align*}
 |\mathbf q|^2
 &=(a_2b_3-a_3b_2)^2+(a_3b_1-a_1b_3)^2+(a_1b_2-a_2b_1)^2\\
 &=(a_1^2+a_2^2+a_3^2)(b_1^2+b_2^2+b_3^2)
   -(a_1b_1+a_2b_2+a_3b_3)^2\\
 &=|\mathbf a|^2|\mathbf b|^2-(\mathbf a\cdot\mathbf b)^2\\
 &=|\mathbf a|^2|\mathbf b|^2\sin^2\theta.
\end{align*}
The magnitude condition on $\mathbf c=t\mathbf q$ gives $t=\pm1$.
The right-handed convention chooses $t=1$, consistently with
$\mathbf i\times\mathbf j=\mathbf k$.
Therefore
\begin{align*}
 \mathbf a\times\mathbf b
 &=(a_2b_3-a_3b_2)\mathbf i
 +(a_3b_1-a_1b_3)\mathbf j
 +(a_1b_2-a_2b_1)\mathbf k.
\end{align*}

\subsection{Properties and areas}
The definition and component formula imply
\begin{align*}
 \mathbf a\cdot(\mathbf a\times\mathbf b)&=0,\\
 \mathbf a\times\mathbf b&=-\mathbf b\times\mathbf a,\\
 (\mathbf a+\mathbf b)\times\mathbf c
   &=\mathbf a\times\mathbf c+\mathbf b\times\mathbf c,\\
 |\mathbf a\times\mathbf b|^2+(\mathbf a\cdot\mathbf b)^2
   &=|\mathbf a|^2|\mathbf b|^2.
\end{align*}
In particular,
\[
 \mathbf i\times\mathbf j=\mathbf k,\quad
 \mathbf j\times\mathbf k=\mathbf i,\quad
 \mathbf k\times\mathbf i=\mathbf j,\qquad
 \mathbf i\times\mathbf i=\mathbf j\times\mathbf j=\mathbf k\times\mathbf k=\mathbf0.
\]
A parallelogram with adjacent edges $\mathbf a,\mathbf b$ has height
$|\mathbf b|\sin\theta$ above a base of length $|\mathbf a|$, so
\[
 \operatorname{Area}(\text{parallelogram})=|\mathbf a\times\mathbf b|,
 \qquad
 \operatorname{Area}(\text{triangle})=\frac12|\mathbf a\times\mathbf b|.
\]

\section{Triple products}

The \emph{scalar triple product} is
\[
 [\mathbf a,\mathbf b,\mathbf c]=\mathbf a\cdot(\mathbf b\times\mathbf c).
\]
For a parallelepiped with edges $\mathbf a,\mathbf b,\mathbf c$, take the
base to be the parallelogram spanned by $\mathbf b,\mathbf c$.
For a nondegenerate base its area is $|\mathbf b\times\mathbf c|$.
The height is the magnitude of the projection of $\mathbf a$ onto the normal:
\begin{align*}
 h
 &=\left|\frac{\mathbf a\cdot(\mathbf b\times\mathbf c)}
 {|\mathbf b\times\mathbf c|^2}(\mathbf b\times\mathbf c)\right|\\
 &=\frac{|\mathbf a\cdot(\mathbf b\times\mathbf c)|}{|\mathbf b\times\mathbf c|}.
\end{align*}
Hence the volume is
\[
 V=|\mathbf b\times\mathbf c|h
 =\boxed{|[\mathbf a,\mathbf b,\mathbf c]|}.
\]
This volume formula also gives zero for a degenerate parallelepiped.
The scalar triple product is signed volume: it is positive for a right-handed
triple and negative for a left-handed triple. It is unchanged by cyclic
permutations and changes sign when any pair is swapped:
\begin{align*}
 [\mathbf a,\mathbf b,\mathbf c]
 &=[\mathbf b,\mathbf c,\mathbf a]=[\mathbf c,\mathbf a,\mathbf b],\\
 [\mathbf a,\mathbf b,\mathbf c]&=-[\mathbf b,\mathbf a,\mathbf c].
\end{align*}
The \emph{vector triple product} $\mathbf a\times(\mathbf b\times\mathbf c)$
will be evaluated using index notation below.

\section{Index notation}

Relabel the standard basis as
$\mathbf e_1=\mathbf i$, $\mathbf e_2=\mathbf j$, $\mathbf e_3=\mathbf k$.
\textbf{Einstein notation} convention suppresses the summation sign for a repeated
index: here every summed index runs from $1$ to $3$. For example,
\[
 \mathbf a=a_i\mathbf e_i
 =\sum_{i=1}^3a_i\mathbf e_i.
\]
A free index is not summed; a repeated dummy index can be renamed consistently.

\subsection{The Kronecker delta and the dot product}
Define
\[
 \delta_{ij}=\mathbf e_i\cdot\mathbf e_j
 =\begin{cases}1,&i=j,\\0,&i\ne j.\end{cases}
\]
Then
\begin{align*}
 \mathbf a\cdot\mathbf b
 &=(a_i\mathbf e_i)\cdot(b_j\mathbf e_j)\\
 &=a_i b_j(\mathbf e_i\cdot\mathbf e_j)
 =\delta_{ij}a_i b_j=a_i b_i.
\end{align*}
The useful contraction rules are
$\delta_{ij}A_j=A_i$ and $\delta_{ij}A_i=A_j$.

\subsection{The Levi--Civita symbol and the cross product}
Expand a cross product of basis vectors in the same basis:
\[
 \mathbf e_i\times\mathbf e_j=\varepsilon_{ijk}\mathbf e_k,
 \qquad
 \varepsilon_{ijk}=(\mathbf e_i\times\mathbf e_j)\cdot\mathbf e_k
 =[\mathbf e_i,\mathbf e_j,\mathbf e_k].
\]
The last equality follows from the cyclic property of the scalar triple product.
Since the basis is right-handed and orthonormal, $\varepsilon_{123}=1$ and
\[
 \varepsilon_{ijk}=
 \begin{cases}
 +1,&(i,j,k)=(1,2,3),(2,3,1),(3,1,2),\\
 -1,&(i,j,k)=(1,3,2),(3,2,1),(2,1,3),\\
 0,&\text{any two indices coincide}.
 \end{cases}
\]
Thus $\varepsilon_{ijk}=\varepsilon_{jki}=\varepsilon_{kij}$ and swapping
any two indices changes its sign. The cross product becomes
\[
 \boxed{\mathbf a\times\mathbf b
 =a_i b_j(\mathbf e_i\times\mathbf e_j)
 =\varepsilon_{ijk}a_i b_j\mathbf e_k}.
\]

\subsection{Contracting two Levi--Civita symbols}
The identity needed for a vector triple product is
\[
 \boxed{\varepsilon_{kij}\varepsilon_{kn\ell}
 =\delta_{in}\delta_{j\ell}-\delta_{i\ell}\delta_{jn}}.
\]
There is a sum over $k$. A nonzero term requires $i\ne j$, $n\ne\ell$
and a $k$ different from all four indices. Consequently there are two
nonzero cases: $i=n\ne j=\ell$ gives $+1$, while
$i=\ell\ne j=n$ gives $-1$. In every other case the sum is zero,
which is exactly the right-hand side.

\subsection{The vector triple product}
\begin{align*}
 \mathbf a\times(\mathbf b\times\mathbf c)
 &=\mathbf a\times(\varepsilon_{ijk}b_i c_j\mathbf e_k)\\
 &=\varepsilon_{ijk}b_i c_j a_\ell
       (\mathbf e_\ell\times\mathbf e_k)\\
 &=\varepsilon_{ijk}\varepsilon_{\ell kn}a_\ell b_i c_j\mathbf e_n\\
 &=\varepsilon_{kij}\varepsilon_{kn\ell}a_\ell b_i c_j\mathbf e_n\\
 &=(\delta_{in}\delta_{j\ell}-\delta_{i\ell}\delta_{jn})
       a_\ell b_i c_j\mathbf e_n\\
 &=(a_jc_j)b_n\mathbf e_n-(a_i b_i)c_n\mathbf e_n\\
 &=(\mathbf a\cdot\mathbf c)\mathbf b-(\mathbf a\cdot\mathbf b)\mathbf c.
\end{align*}
Therefore
\[
 \boxed{\mathbf a\times(\mathbf b\times\mathbf c)
 =(\mathbf a\cdot\mathbf c)\mathbf b-(\mathbf a\cdot\mathbf b)\mathbf c}.
\]
Week 2 ends with this calculation unfinished; Week 3 supplies the contraction
identity and completes it. The recap is combined here rather than duplicated.

\section{A determinant mnemonic and a plane example}

The component formula can be remembered by expanding formally along the first row:
\begin{align*}
 \mathbf a\times\mathbf b
 &=\begin{vmatrix}
 \mathbf i&\mathbf j&\mathbf k\\a_1&a_2&a_3\\b_1&b_2&b_3
 \end{vmatrix}\\
 &=\mathbf i\begin{vmatrix}a_2&a_3\\b_2&b_3\end{vmatrix}
 -\mathbf j\begin{vmatrix}a_1&a_3\\b_1&b_3\end{vmatrix}
 +\mathbf k\begin{vmatrix}a_1&a_2\\b_1&b_2\end{vmatrix},
\end{align*}
where $\left|\begin{smallmatrix}a&b\\c&d\end{smallmatrix}\right|=ad-bc$.
The vector-valued first row is a mnemonic for cofactor expansion.

\begin{example}[Lecture Example 5: plane through three points and distance]
Let
\[
 A=(-1,0,1),\qquad B=(2,-2,1),\qquad C=(3,-3,0),\qquad D=(-3,-1,2).
\]
Find a Cartesian equation of the plane through $A,B,C$ and its distance from $D$.
A normal vector is
\begin{align*}
 \mathbf n
 &=\overrightarrow{AB}\times\overrightarrow{AC}
 =\langle3,-2,0\rangle\times\langle4,-3,-1\rangle\\
 &=\mathbf i\bigl((-2)(-1)-0(-3)\bigr)
 -\mathbf j\bigl(3(-1)-0(4)\bigr)
 +\mathbf k\bigl(3(-3)-(-2)(4)\bigr)\\
 &=2\mathbf i+3\mathbf j-\mathbf k=\langle2,3,-1\rangle.
\end{align*}
With $\mathbf r_A=\langle-1,0,1\rangle$,
$\mathbf r_A\cdot\mathbf n=-3$. Thus the plane is
\[
 \boxed{2x+3y-z=-3}.
\]
The shortest distance is the length of the projection of $\overrightarrow{AD}$
onto the normal. Since $\overrightarrow{AD}=\langle-2,-1,1\rangle$,
\begin{align*}
 d
 &=\left|\operatorname{proj}_{\mathbf n}\overrightarrow{AD}\right|
 =\frac{|\overrightarrow{AD}\cdot\mathbf n|}{|\mathbf n|}\\
 &=\frac{|-4-3-1|}{\sqrt{4+9+1}}
 =\frac8{\sqrt{14}}=\boxed{\frac{4\sqrt{14}}7}.
\end{align*}
\end{example}


